\documentclass{article}
\usepackage[utf8]{inputenc}
\usepackage[T1]{fontenc}
\usepackage{geometry}
\usepackage{graphicx}
\usepackage{float}
\usepackage{hyperref}
\usepackage[english,indonesian]{babel}

\title{Optimasi Pemrosesan Template LaTeX Menggunakan Pendekatan Berbasis AI}
\author{Academic Support Agent}
\date{\today}

\begin{document}
\maketitle

\begin{abstract}

\end{abstract}

\chapter{PENDAHULUAN}


\chapter{TINJAUAN PUSTAKA}


% Original content (commented out):
% \documentclass[12pt]{article}

\input{1_header/packages.tex}

\begin{document}
\input{1_header/title.tex}
\restoregeometry

\section*{Acknowledgments}
\input{1_header/acknowledgments.tex}
\pagebreak

\setcounter{tocdepth}{5}
\tableofcontents
\pagebreak

\listoffigures
\listoftables
\pagebreak

\justify

\section*{Abstract}
\input{1_header/abstract.tex}
\pagebreak

\section{Introduction}
\input{1_header/introduction.tex}

\section{Literature Review}
\input{2_body/literature.tex}

\section{Methods}
\input{2_body/methods.tex}

\section{Results}
\input{2_body/results.tex}

\section{Discussion}
\input{2_body/discussion.tex}

\section{Conclusions}
\input{2_body/conclusions.tex}

\section{Recommendations}
\input{2_body/recommendations.tex}

\section{References}
\label{sec:ref}
\printbibliography[heading=none]

\section{Appendices}
\input{3_footer/appendices.tex}

\end{document}


\end{document}